\begin{table}[H]
\centering
\scriptsize
\caption{Dodatkowy miesięczny koszt BDP i jego ekwiwalent na osobę}
\label{tab:per-capita-equivalents}
\begin{tabular}{@{}rrrrr@{}}
\toprule
BDP w modelu & Dodatkowe transfery & Ekwiwalent & Koszt brutto & \(\Delta\) deficytu \\
PLN/HH/mies. & mld PLN/mies. & PLN/os./mies. & \% PKB bez BDP, m60 & pp PKB \\
\midrule
500  & 10,8 & 283   & 2,40  & 1,50 \\
1000 & 21,5 & 565   & 4,80  & 3,25 \\
1500 & 32,2 & 848   & 7,20  & 4,79 \\
2000 & 43,0 & 1 130 & 9,59  & 6,51 \\
2500 & 53,8 & 1 415 & 12,00 & 7,76 \\
3000 & 64,6 & 1 700 & 14,42 & 8,86 \\
\bottomrule
\end{tabular}
\vspace{0.25em}

\parbox{\textwidth}{\scriptsize \textit{Uwaga:} BDP w modelu oznacza miesięczny transfer na agenta reprezentującego gospodarstwo domowe, a nie kwotę na mieszkańca Polski. Dodatkowe transfery są liczone jako różnica modelowych miesięcznych wydatków na transfery społeczne względem wariantu bez BDP. Przykładowo dla BDP 500~PLN całkowite transfery społeczne wynoszą 16,77 mld PLN/mies., z czego 6,02 mld PLN/mies. przypada na transfery istniejące w wariancie bez BDP; różnica daje 10,8 mld PLN/mies. Ekwiwalent na osobę to ta różnica podzielona przez populację referencyjną 38 mln osób. Koszt brutto wyrażono względem miesięcznego PKB wariantu bez BDP po 60 miesiącach, aby wszystkie poziomy transferu miały ten sam mianownik.}
\end{table}
